% Resumen del mundo «Fábrica de Funciones» (Cálculo Diferencial, mundo 2).
% Lecciones alg-52, calc-tipos, alg-53, calc-tramos, alg-61, alg-62. Usa los ejemplos de los videos; el taller trae ejercicios nuevos.
\documentclass[11pt,a4paper]{article}
\usepackage{zuka-resumen}
\usepackage{needspace}
\zukacurso{calc}
\zukamundo{2}{Fábrica de Funciones}{Funciones}

% Cada nivel empieza con espacio suficiente (el título no queda solo al pie de la página).
\let\nivelorig\nivel
\renewcommand{\nivel}[2]{\Needspace{8\baselineskip}\nivelorig{#1}{#2}}
% Dentro de una minipage las cajas no se parten entre páginas.
\AtBeginEnvironment{minipage}{\tcbset{zbase/.append style={unbreakable}}}
% Panel pequeño de gráfica con ejes en el centro y misma escala en x y en y.
\pgfplotsset{mini/.style={zplot,x=0.42cm,y=0.42cm,xtick=\empty,ytick=\empty,grid=none}}
\tikzset{curvabase/.style={zmute,dashed,line width=0.9pt}}

\begin{document}

\portada{%
Las funciones son el idioma del cálculo: describen cómo una cantidad depende de otra. En este mundo aprendiste a
reconocer una \clave{función}, a hallar su \clave{dominio} y su \clave{rango}, a distinguir sus tipos, a operarlas y
componerlas, a trabajar con funciones \clave{a tramos}, a mover sus gráficas y a deshacerlas con la \clave{inversa}.
Este resumen reúne las definiciones, la notación y los ejemplos de los videos, y termina con un taller.}{%
\item \textbf{Nivel 2-1.} Funciones: dominio, rango, simetría (par e impar) y monotonía.
\item \textbf{Nivel 2-2.} Tipos de funciones: constante, lineal, cuadrática, polinómica, potencia y racional.
\item \textbf{Nivel 2-3.} Operaciones entre funciones y composición.
\item \textbf{Nivel 2-4.} Funciones a tramos, valor absoluto y funciones algebraicas.
\item \textbf{Nivel 2-5.} Transformaciones: trasladar, estirar y reflejar gráficas.
\item \textbf{Nivel 2-6.} Funciones inyectivas, biyectivas e inversas.}

\begin{clavebox}[LAS IDEAS QUE SOSTIENEN TODO EL MUNDO]
\renewcommand{\arraystretch}{1.5}\begin{tabular}{@{}>{\centering\arraybackslash}p{46mm}@{\quad}>{\centering\arraybackslash}p{50mm}@{\quad}>{\centering\arraybackslash}p{50mm}@{}}
$x\ \longmapsto\ f(x)$ & $(g\circ f)(x)=g\big(f(x)\big)$ & $f^{-1}\big(f(x)\big)=x$\\[2pt]
\small cada entrada tiene exactamente una salida & \small primero actúa $f$, luego $g$ & \small la inversa deshace lo que hizo $f$
\end{tabular}
\end{clavebox}

%=====================================================================
\nivel{2-1}{Funciones: dominio, rango y simetría}

\begin{definicion}[Función]
Una \clave{función} $f$ de un conjunto $A$ en un conjunto $B$, $f\colon A\to B$, es una regla que asigna a cada
$x\in A$ \emph{exactamente un} elemento $y=f(x)\in B$. A $x$ se le llama variable independiente (entrada) y a $y$,
variable dependiente (salida).
\end{definicion}

Es como una máquina: entra un número y sale uno solo. La \clave{gráfica} de $f$ es el conjunto de puntos
$\{(x,f(x))\tq x\in A\}$. Una curva del plano es gráfica de una función si y solo si \emph{ninguna recta vertical la
corta dos veces} (prueba de la recta vertical).

\begin{definicion}[Dominio y rango]
El \clave{dominio} $\operatorname{Dom} f$ es el conjunto de entradas permitidas; el \clave{rango}
$\operatorname{Ran} f=\{f(x)\tq x\in\operatorname{Dom} f\}$ es el conjunto de salidas. Si solo se da la fórmula, el
\clave{dominio natural} son todos los reales para los que la fórmula tiene sentido.
\end{definicion}

\begin{minipage}[c]{0.55\linewidth}
Para el dominio natural se excluye lo que no se puede hacer en $\R$:
\begin{itemize}
\item dividir entre cero: $\dfrac{1}{x-2}$ tiene dominio $\R\setminus\{2\}$;
\item sacar raíz par de un negativo: $\sqrt{x+1}$ pide $x+1\ge 0$, así que su dominio es $[-1,\infty)$.
\end{itemize}
\smallskip
Geométricamente, el dominio es la \emph{sombra} de la gráfica sobre el eje $x$ y el rango, su sombra sobre el eje $y$.
Para $f(x)=x^2-1$ el dominio es $\R$ y, como $x^2\ge 0$, se tiene $x^2-1\ge-1$: el rango es $[-1,\infty)$.
\end{minipage}\hfill
\begin{minipage}[c]{0.4\linewidth}\centering
\begin{tikzpicture}
\begin{axis}[zplot,x=0.62cm,y=0.62cm,xmin=-3,xmax=3,ymin=-2,ymax=5.5,xtick={-2,...,2},ytick={-1,...,5}]
  \addplot[zaccent,domain=-2.55:2.55,samples=60] {x^2-1};
  \draw[zcoral,line width=3pt,opacity=0.7] (axis cs:0,-1)--(axis cs:0,5.4);
  \fill[zcoral] (axis cs:0,-1) circle (2.4pt);
  \node[font=\footnotesize,text=zcoral,anchor=west] at (axis cs:0.15,4.6) {rango};
  \node[font=\footnotesize,text=zaccent,anchor=west] at (axis cs:1.35,-1.4) {$y=x^2-1$};
\end{axis}
\end{tikzpicture}
\end{minipage}

\subsection{Simetría y monotonía}
\begin{definicion}[Funciones par e impar]
$f$ es \clave{par} si $f(-x)=f(x)$ para todo $x$ de su dominio: su gráfica es simétrica respecto al eje $y$.
$f$ es \clave{impar} si $f(-x)=-f(x)$: su gráfica es simétrica respecto al origen (media vuelta).
\end{definicion}

\begin{figura}[Izquierda: $x^2$ es par (espejo en el eje $y$). Derecha: $x^3$ es impar (media vuelta en el origen).]
\begin{tikzpicture}
\begin{axis}[mini,x=0.55cm,y=0.55cm,xmin=-4.2,xmax=4.2,ymin=-0.8,ymax=4.6]
  \addplot[zaccent,domain=-2.1:2.1,samples=50] {x^2};
  \draw[zmute,densely dotted] (axis cs:-1.5,2.25)--(axis cs:1.5,2.25);
  \fill[zcoral] (axis cs:-1.5,2.25) circle (2.2pt); \fill[zcoral] (axis cs:1.5,2.25) circle (2.2pt);
  \node[font=\scriptsize,anchor=west] at (axis cs:1.65,2.0) {$(a,f(a))$};
  \node[font=\scriptsize,anchor=east] at (axis cs:-1.65,2.0) {$(-a,f(a))$};
\end{axis}
\end{tikzpicture}\qquad\qquad
\begin{tikzpicture}
\begin{axis}[mini,x=0.55cm,y=0.55cm,xmin=-4.2,xmax=4.2,ymin=-2.7,ymax=2.7]
  \addplot[zaccent,domain=-1.38:1.38,samples=50] {x^3};
  \draw[zmute,densely dotted] (axis cs:-1.2,-1.728)--(axis cs:1.2,1.728);
  \fill[zcoral] (axis cs:-1.2,-1.728) circle (2.2pt); \fill[zcoral] (axis cs:1.2,1.728) circle (2.2pt);
  \node[font=\scriptsize,anchor=west] at (axis cs:1.3,1.6) {$(a,f(a))$};
  \node[font=\scriptsize,anchor=east] at (axis cs:-1.3,-1.6) {$(-a,-f(a))$};
\end{axis}
\end{tikzpicture}
\end{figura}

\begin{definicion}[Creciente y decreciente]
$f$ es \clave{creciente} en un intervalo $I$ si $x_1<x_2$ en $I$ implica $f(x_1)<f(x_2)$: al avanzar a la derecha la
gráfica sube. Es \clave{decreciente} si $x_1<x_2$ implica $f(x_1)>f(x_2)$.
\end{definicion}
\begin{ejemplo}[Del video]
$f(x)=x^2$ decrece en $(-\infty,0]$ y crece en $[0,\infty)$: la parábola baja hasta el vértice y luego sube.
\end{ejemplo}

\begin{ojo}[Evaluar con paréntesis]
Para calcular $f(-3)$ con $f(x)=x^2+2x$ se reemplaza cada $x$ por $(-3)$: $(-3)^2+2(-3)=9-6=3$. Sin paréntesis
se cae en el error $-3^2=-9$.
\end{ojo}

%=====================================================================
\nivel{2-2}{Tipos de funciones}

\begin{center}
\begin{tabularx}{\linewidth}{@{}l l L@{}}
\toprule
Tipo & Fórmula & Gráfica y rasgos\\ \midrule
Constante & $f(x)=c$ & recta horizontal; rango $\{c\}$\\
Lineal & $f(x)=mx+b$, $m\ne 0$ & recta de pendiente $m$ y corte $b$\\
Cuadrática & $f(x)=ax^2+bx+c$, $a\ne0$ & parábola; vértice en $x=-\frac{b}{2a}$\\
Polinómica & $f(x)=a_nx^n+\dots+a_1x+a_0$ & dominio $\R$; a lo sumo $n$ raíces reales\\
Potencia & $f(x)=x^{a}$, $a$ constante & $x^2,\ x^3,\ x^{1/2}=\sqrt{x},\ x^{-1}=\frac1x$\\
Racional & $f(x)=\dfrac{P(x)}{Q(x)}$, $P,Q$ polinomios & dominio: $Q(x)\ne0$; posibles asíntotas\\
\bottomrule
\end{tabularx}
\end{center}

\begin{propiedad}[Raíces y factores de un polinomio]
$r$ es raíz de $P$ (es decir, $P(r)=0$) si y solo si $(x-r)$ es un factor de $P(x)$. Si el factor aparece un número
impar de veces, la gráfica \emph{cruza} el eje $x$ en $r$; si aparece un número par de veces, lo \emph{toca y rebota}.
\end{propiedad}

\begin{minipage}[t]{0.48\linewidth}
\begin{ejemplo}[Del video]
$f(x)=(x+1)(x-2)^2$ corta el eje en $-1$ (raíz simple) y rebota en $2$ (raíz doble). Su corte con el eje $y$ es
$f(0)=1\cdot4=4$. El coeficiente principal es positivo, así que a la derecha la gráfica sube.
\end{ejemplo}
\end{minipage}\hfill
\begin{minipage}[t]{0.48\linewidth}
\begin{ejemplo}[Del video]
$f(x)=\dfrac{1}{x-2}+1$ no existe en $x=2$: allí hay una \clave{asíntota vertical}. Para $x$ muy grande,
$\frac{1}{x-2}\approx0$ y $f(x)\approx1$: la recta $y=1$ es una \clave{asíntota horizontal}.
\end{ejemplo}
\end{minipage}

\begin{figura}[Izquierda: raíz simple en $-1$ y raíz doble en $2$. Derecha: asíntotas $x=2$ y $y=1$.]
\begin{tikzpicture}
\begin{axis}[zplot,x=0.62cm,y=0.42cm,xmin=-2.2,xmax=3.8,ymin=-4.3,ymax=7,xtick={-1,1,2,3},ytick={-4,-2,2,4,6}]
  \addplot[zaccent,domain=-1.35:3.25,samples=80] {(x+1)*(x-2)^2};
  \fill[zcoral] (axis cs:-1,0) circle (2.2pt); \fill[zcoral] (axis cs:2,0) circle (2.2pt);
\end{axis}
\end{tikzpicture}\qquad
\begin{tikzpicture}
\begin{axis}[zplot,x=0.5cm,y=0.5cm,xmin=-2.5,xmax=6.5,ymin=-3.3,ymax=5.5,xtick={-2,2,4,6},ytick={-2,1,3,5}]
  \addplot[curvabase] coordinates {(2,-3.3) (2,5.5)};
  \addplot[curvabase] coordinates {(-2.5,1) (6.5,1)};
  \addplot[zaccent,domain=-2.4:1.75,samples=60] {1/(x-2)+1};
  \addplot[zaccent,domain=2.23:6.4,samples=60] {1/(x-2)+1};
  \node[font=\scriptsize,text=zmute,anchor=south west] at (axis cs:2.1,4.6) {$x=2$};
  \node[font=\scriptsize,text=zmute,anchor=south] at (axis cs:5.4,1.05) {$y=1$};
\end{axis}
\end{tikzpicture}
\end{figura}

%=====================================================================
\Needspace{16\baselineskip}\nivel{2-3}{Operaciones y composición de funciones}

\begin{definicion}[Operaciones punto a punto]
Dadas $f$ y $g$, en cada $x$ se operan las alturas:
\[ (f\pm g)(x)=f(x)\pm g(x),\qquad (fg)(x)=f(x)\,g(x),\qquad \left(\frac{f}{g}\right)(x)=\frac{f(x)}{g(x)}. \]
El dominio es $\operatorname{Dom}f\cap\operatorname{Dom}g$, quitando además los $x$ con $g(x)=0$ en el cociente.
\end{definicion}

\begin{ejemplo}
Con $f(x)=\sqrt{x}$ y $g(x)=x-3$: $(f+g)(x)=\sqrt{x}+x-3$ tiene dominio $[0,\infty)$, mientras que
$\left(\frac fg\right)(x)=\dfrac{\sqrt{x}}{x-3}$ tiene dominio $[0,3)\cup(3,\infty)$.
\end{ejemplo}

\begin{definicion}[Composición]
La \clave{composición} de $g$ con $f$ es $(g\circ f)(x)=g\big(f(x)\big)$: primero actúa $f$ y su salida entra a $g$.
Su dominio son los $x\in\operatorname{Dom} f$ tales que $f(x)\in\operatorname{Dom} g$.
\end{definicion}

\begin{figura}[La composición encadena máquinas. Con $f(x)=x+1$ y $g(x)=x^2$: $g(f(3))=g(4)=16$.]
\begin{tikzpicture}[mq/.style={draw=zaccent,fill=zaccent!10,rounded corners=3pt,minimum width=20mm,minimum height=10mm,font=\small,line width=0.9pt},font=\small]
  \node (x) at (0,0) {$3$};
  \node[mq] (f) at (2.4,0) {$f:\ \square+1$};
  \node[mq] (g) at (6.4,0) {$g:\ (\square)^2$};
  \node (y) at (8.9,0) {$16$};
  \draw[->,line width=0.9pt] (x)--(f); \draw[->,line width=0.9pt] (f)--(g) node[midway,above,font=\footnotesize] {$f(3)=4$};
  \draw[->,line width=0.9pt] (g)--(y);
  \node[font=\footnotesize,text=zmute] at (4.4,-1) {$(g\circ f)(x)=(x+1)^2$};
\end{tikzpicture}
\end{figura}

\begin{ejemplo}[Del video: el orden importa]
Con $f(x)=x+1$ y $g(x)=x^2$:
\[ (g\circ f)(3)=g(4)=16,\qquad (f\circ g)(3)=f(9)=10. \]
En general, $(g\circ f)(x)=(x+1)^2$ y $(f\circ g)(x)=x^2+1$: la composición \emph{no} es conmutativa, $g\circ f\ne f\circ g$.
\end{ejemplo}

\begin{ejemplo}[Del video: desarmar una compuesta]
En $h(x)=(2x+1)^3$ la operación de adentro es $u=2x+1$ y la de afuera es $u^3$. Así $h=G\circ F$ con
$F(x)=2x+1$ y $G(u)=u^3$. Esta idea será clave para la regla de la cadena.
\end{ejemplo}

\begin{ojo}[$(f\circ g)(x)$ no es $f(x)\cdot g(x)$]
El círculo $\circ$ significa «aplicar una dentro de la otra», no multiplicar.
\end{ojo}

%=====================================================================
\nivel{2-4}{Funciones a tramos y valor absoluto}

\begin{definicion}[Función a tramos]
Una función \clave{a tramos} usa una fórmula distinta en cada parte de su dominio. Para evaluarla, primero se mira
en qué tramo cae $x$ y luego se usa \emph{solo} la fórmula de ese tramo.
\end{definicion}

\begin{ejemplo}[Un parqueadero]
Un parqueadero cobra \$3\,000 por la primera hora y \$2\,000 por cada hora adicional (proporcional al tiempo):
\[ P(t)=\begin{cases} 3\,000, & 0<t\le 1,\\ 3\,000+2\,000\,(t-1), & t>1.\end{cases} \]
Para $t=2.5$ h se usa el segundo tramo: $P(2.5)=3\,000+2\,000(1.5)=6\,000$ pesos. En la gráfica, el punto del
corte $t=1$ va lleno en el tramo cuya condición trae $\le$ y vacío en el otro.
\end{ejemplo}

\begin{propiedad}[El valor absoluto es una función a tramos]
\[ \abs{x}=\begin{cases} x, & x\ge0,\\ -x, & x<0.\end{cases} \]
Su gráfica es una V con vértice en el origen. Para graficar $\abs{g(x)}$ se deja igual lo que $g$ tiene sobre el
eje $x$ y se \emph{dobla hacia arriba} lo que tiene debajo.
\end{propiedad}

\begin{minipage}[c]{0.5\linewidth}
\begin{ejemplo}[Del video: $\abs{x^2-4}$]
La parábola $x^2-4$ está bajo el eje entre sus raíces $-2$ y $2$. Ese trozo se refleja hacia arriba:
\[ \abs{x^2-4}=\begin{cases} x^2-4, & \abs{x}\ge 2,\\ 4-x^2, & \abs{x}<2.\end{cases} \]
\end{ejemplo}
Una función \clave{algebraica} se arma con sumas, productos, cocientes y raíces de $x$: por ejemplo
$\sqrt{x}$, $\dfrac{x}{x+1}$ o $x\sqrt{x+3}$.
\end{minipage}\hfill
\begin{minipage}[c]{0.45\linewidth}\centering
\begin{tikzpicture}
\begin{axis}[zplot,x=0.7cm,y=0.42cm,xmin=-3.2,xmax=3.2,ymin=-4.8,ymax=5.8,xtick={-2,...,2},ytick={-4,-2,2,4}]
  \addplot[curvabase,domain=-3:3,samples=60] {x^2-4};
  \addplot[zaccent,domain=-3:3,samples=121] {abs(x^2-4)};
\end{axis}
\end{tikzpicture}
\end{minipage}

%=====================================================================
\Needspace{16\baselineskip}\nivel{2-5}{Transformaciones de funciones}

\begin{propiedad}[Cómo se mueve la gráfica de $y=f(x)$, con $h,k>0$ y $a>1$]
\renewcommand{\arraystretch}{1.3}
\begin{tabularx}{\linewidth}{@{}l L l L@{}}
$y=f(x)+k$ & sube $k$ & $y=f(x)-k$ & baja $k$\\
$y=f(x-h)$ & corre $h$ a la \emph{derecha} & $y=f(x+h)$ & corre $h$ a la \emph{izquierda}\\
$y=a\,f(x)$ & estira en vertical (factor $a$) & $y=\tfrac1a f(x)$ & encoge en vertical\\
$y=f(ax)$ & encoge en horizontal (factor $\tfrac1a$) & $y=f\!\left(\tfrac xa\right)$ & estira en horizontal\\
$y=-f(x)$ & refleja en el eje $x$ & $y=f(-x)$ & refleja en el eje $y$
\end{tabularx}
\end{propiedad}

Regla para recordar: \emph{lo de afuera} actúa en vertical y tal como se lee; \emph{lo de adentro}, junto a la $x$,
actúa en horizontal y al revés de lo que parece.

\begin{figura}[Las transformaciones aplicadas a $y=x^2$ (punteada).]
\begin{tikzpicture}
\begin{axis}[mini,x=0.36cm,y=0.36cm,xmin=-3.2,xmax=3.2,ymin=-1,ymax=6.5,title={\footnotesize $y=x^2+2$}]
  \addplot[curvabase,domain=-2.5:2.5,samples=40] {x^2};
  \addplot[zaccent,domain=-2.1:2.1,samples=40] {x^2+2};
\end{axis}
\end{tikzpicture}\enspace
\begin{tikzpicture}
\begin{axis}[mini,x=0.36cm,y=0.36cm,xmin=-2.2,xmax=5.6,ymin=-1,ymax=6.5,title={\footnotesize $y=(x-3)^2$}]
  \addplot[curvabase,domain=-2.1:2.5,samples=40] {x^2};
  \addplot[zaccent,domain=0.5:5.5,samples=40] {(x-3)^2};
\end{axis}
\end{tikzpicture}\enspace
\begin{tikzpicture}
\begin{axis}[mini,x=0.36cm,y=0.36cm,xmin=-3.2,xmax=3.2,ymin=-1,ymax=6.5,title={\footnotesize $y=2x^2$}]
  \addplot[curvabase,domain=-2.5:2.5,samples=40] {x^2};
  \addplot[zaccent,domain=-1.78:1.78,samples=40] {2*x^2};
\end{axis}
\end{tikzpicture}\enspace
\begin{tikzpicture}
\begin{axis}[mini,x=0.36cm,y=0.36cm,xmin=-3.2,xmax=3.2,ymin=-4,ymax=4,title={\footnotesize $y=-x^2$}]
  \addplot[curvabase,domain=-1.95:1.95,samples=40] {x^2};
  \addplot[zaccent,domain=-1.95:1.95,samples=40] {-x^2};
\end{axis}
\end{tikzpicture}
\end{figura}

\begin{ejemplo}[Del video: el orden importa]
Partiendo de $\sqrt{x}$: estirar $\times2$ y luego subir $1$ da $2\sqrt{x}+1$; subir $1$ y luego estirar da
$2(\sqrt{x}+1)=2\sqrt{x}+2$. Son funciones distintas: el estiramiento posterior también duplica la subida.
\end{ejemplo}

%=====================================================================
\nivel{2-6}{Función biyectiva e inversa}

\begin{definicion}[Inyectiva, sobreyectiva y biyectiva]
$f\colon A\to B$ es \clave{inyectiva} si $x_1\ne x_2$ implica $f(x_1)\ne f(x_2)$: ninguna recta horizontal corta su
gráfica dos veces. Es \clave{sobreyectiva} si su rango es todo $B$. Es \clave{biyectiva} si es inyectiva y sobreyectiva.
\end{definicion}
$x^3$ es inyectiva en $\R$; $x^2$ no lo es, porque $f(2)=f(-2)=4$.

\begin{definicion}[Función inversa]
Si $f\colon A\to B$ es biyectiva, su \clave{inversa} $f^{-1}\colon B\to A$ está definida por
\[ f^{-1}(y)=x \iff f(x)=y. \]
Así $f^{-1}\big(f(x)\big)=x$ para $x\in A$ y $f\big(f^{-1}(y)\big)=y$ para $y\in B$. El dominio de $f^{-1}$ es el rango de $f$.
\end{definicion}

La inversa deshace en orden contrario, como al quitarse los zapatos antes que las medias.

\begin{minipage}[c]{0.52\linewidth}
\begin{ejemplo}[Del video: hallar la inversa]
Para $f(x)=2x+3$ se escribe $y=2x+3$ y se despeja $x$:
\[ x=\frac{y-3}{2}\quad\Longrightarrow\quad f^{-1}(x)=\frac{x-3}{2}. \]
\end{ejemplo}
\begin{propiedad}[Gráfica de la inversa]
Las gráficas de $f$ y $f^{-1}$ son reflejos en la recta $y=x$: si $(a,b)$ está en $f$, entonces $(b,a)$ está en $f^{-1}$.
\end{propiedad}
\end{minipage}\hfill
\begin{minipage}[c]{0.44\linewidth}\centering
\begin{tikzpicture}
\begin{axis}[zplot,x=0.4cm,y=0.4cm,xmin=-4.5,xmax=8.5,ymin=-4.5,ymax=8.5,xtick={4,8},ytick={4,8}]
  \addplot[curvabase,domain=-4.5:8.5,samples=2] {x};
  \addplot[zaccent,domain=-3.6:2.6,samples=2] {2*x+3};
  \addplot[zcoral,domain=-4.2:8.2,samples=2] {(x-3)/2};
  \draw[zmute,densely dotted] (axis cs:1,5)--(axis cs:5,1);
  \fill[zaccent] (axis cs:1,5) circle (2.2pt); \fill[zcoral] (axis cs:5,1) circle (2.2pt);
  \node[font=\scriptsize,anchor=north west] at (axis cs:1.3,4.9) {$(1,5)$};
  \node[font=\scriptsize,anchor=north west] at (axis cs:5.2,0.8) {$(5,1)$};
  \node[font=\scriptsize,text=zaccent,anchor=west] at (axis cs:2.7,8.1) {$f$};
  \node[font=\scriptsize,text=zcoral,anchor=south] at (axis cs:7.6,2.4) {$f^{-1}$};
  \node[font=\scriptsize,text=zmute,anchor=north west] at (axis cs:6.4,6.2) {$y=x$};
\end{axis}
\end{tikzpicture}
\end{minipage}

\begin{ejemplo}[Del video: restringir el dominio]
$x^2$ no es inyectiva en $\R$, pero sí en $[0,\infty)$. Con esa restricción su inversa es $\sqrt{x}$:
$\sqrt{x^2}=x$ para $x\ge0$.
\end{ejemplo}

\begin{ojo}[$f^{-1}$ no es $1/f$]
El exponente $-1$ en $f^{-1}$ indica inversa para la composición, no recíproco: para $f(x)=2x+3$,
$f^{-1}(x)=\frac{x-3}{2}$, mientras que $\frac{1}{f(x)}=\frac{1}{2x+3}$.
\end{ojo}

%=====================================================================
\repaso
\begin{center}
\begin{tabularx}{\linewidth}{@{}l L@{}}
\toprule
Idea & Fórmula o regla\\ \midrule
Dominio natural & sin denominadores nulos ni raíces pares de negativos\\
Par / impar & $f(-x)=f(x)$ (eje $y$) \ / \ $f(-x)=-f(x)$ (origen)\\
Raíz y factor & $P(r)=0\iff(x-r)$ divide a $P(x)$; raíz doble: rebota\\
Racional & asíntota vertical donde $Q(x)=0$ (sin cancelarse con $P$)\\
Composición & $(g\circ f)(x)=g(f(x))$; en general $g\circ f\ne f\circ g$\\
Valor absoluto & $\abs{x}=x$ si $x\ge0$, $-x$ si $x<0$; $\abs{g}$ dobla lo negativo hacia arriba\\
Transformaciones & $af(x-h)+k$: corre $h$, estira por $a$, sube $k$; $-f(x)$ y $f(-x)$ reflejan\\
Inversa & $y=f(x)$, despejar $x$, intercambiar; gráfica reflejada en $y=x$\\
\bottomrule
\end{tabularx}
\end{center}

\begin{ojo}[Errores frecuentes]
\begin{itemize}
\item Componer en el orden equivocado: en $f(g(x))$ actúa primero $g$.
\item Correr $f(x+2)$ hacia la derecha (va a la izquierda).
\item Confundir $f^{-1}(x)$ con $\dfrac{1}{f(x)}$.
\end{itemize}
\end{ojo}

%=====================================================================
\taller
\begin{ejercicios}
\item Halla el dominio natural de cada función.
  \begin{ejercicios}
  \item $f(x)=\dfrac{x+1}{x^2-9}$ \qquad (b) $g(x)=\sqrt{4-2x}$ \qquad (c) $h(x)=\dfrac{1}{\sqrt{x+5}}$
  \end{ejercicios}
\item Decide si cada función es par, impar o ninguna de las dos. Justifica con $f(-x)$.
  \begin{ejercicios}
  \item $f(x)=x^4-3x^2$ \qquad (b) $g(x)=\dfrac{x}{x^2+1}$ \qquad (c) $h(x)=x^2+x$
  \end{ejercicios}
\item Grafica $f(x)=3-(x-1)^2$ restringida a $[-1,4)$. Con la gráfica, halla su dominio, su rango y los intervalos
  donde crece y donde decrece.
\item Halla el polinomio de grado $3$ que tiene raíz simple en $x=-2$, raíz doble en $x=1$ y cuya gráfica corta el eje
  $y$ en $4$. Escríbelo desarrollado y describe qué hace la gráfica en $x=1$.
\item Para $f(x)=\dfrac{3x-1}{x+2}$ halla las asíntotas vertical y horizontal y los cortes con los ejes. Haz un bosquejo.
\item Sean $f(x)=\sqrt{x}$ y $g(x)=x^2-4$. Halla cada función y su dominio.
  \begin{ejercicios}
  \item $f\circ g$ \qquad (b) $g\circ f$ \qquad (c) $f\cdot g$
  \end{ejercicios}
\item Escribe cada función como una composición $G\circ F$ de dos funciones más sencillas.
  \begin{ejercicios}
  \item $h(x)=\sqrt{5-x^3}$ \qquad (b) $k(x)=\dfrac{3}{(2x-1)^2}$
  \end{ejercicios}
\item Una empresa de acueducto cobra el agua según el consumo mensual $x$ (en \si{\cubic\metre}):
  \[ C(x)=\begin{cases} 2\,500\,x, & 0\le x\le 20,\\ 50\,000+4\,000\,(x-20), & x>20.\end{cases} \]
  \begin{ejercicios}
  \item Calcula $C(15)$ y $C(28)$.
  \item Dibuja la gráfica. ¿Hay un salto en $x=20$?
  \item Una familia pagó \$70\,000. ¿Cuántos metros cúbicos consumió?
  \end{ejercicios}
\item Transformaciones y valor absoluto.
  \begin{ejercicios}
  \item Describe, en orden, las transformaciones que llevan $y=\sqrt{x}$ a $y=-2\sqrt{x+3}+1$. Da su dominio y su rango,
    y halla a dónde va el punto $(4,2)$ de $y=\sqrt{x}$.
  \item Escribe $f(x)=\abs{3-x}+x$ como una función a tramos sin barras y grafícala.
  \end{ejercicios}
\item Funciones inversas.
  \begin{ejercicios}
  \item Halla la inversa de $f(x)=\dfrac{2x-1}{x+3}$ y su dominio.
  \item Halla la inversa de $g(x)=(x+1)^2$ con $x\ge-1$. Dibuja $g$, $g^{-1}$ y la recta $y=x$ en un mismo plano.
  \end{ejercicios}
\end{ejercicios}

\Needspace{10\baselineskip}\respuestas
\begin{respuestascortas}
\begin{soluciones}
\item (a) $\R\setminus\{-3,3\}$ \ (b) $(-\infty,2]$ \ (c) $(-5,\infty)$.
\item (a) Par. (b) Impar: $g(-x)=\frac{-x}{x^2+1}=-g(x)$. (c) Ninguna: $h(-1)=0$ y $h(1)=2$.
\item Vértice $(1,3)$; $f(-1)=-1$, $f(4)=-6$ (vacío). Dominio $[-1,4)$, rango $(-6,3]$; crece en $[-1,1]$, decrece en $[1,4)$.
\item $f(x)=a(x+2)(x-1)^2$ con $f(0)=2a=4$: $f(x)=2(x+2)(x-1)^2=2x^3-6x+4$. En $x=1$ toca el eje y rebota.
\item AV $x=-2$; AH $y=3$; cortes $\left(\tfrac13,0\right)$ y $\left(0,-\tfrac12\right)$.
\item (a) $\sqrt{x^2-4}$, $(-\infty,-2]\cup[2,\infty)$. (b) $x-4$, pero con dominio $[0,\infty)$. (c) $\sqrt{x}\,(x^2-4)$, $[0,\infty)$.
\item (a) $F(x)=5-x^3$, $G(u)=\sqrt{u}$. (b) $F(x)=2x-1$, $G(u)=\frac{3}{u^2}$. (Hay otras respuestas válidas.)
\item (a) $C(15)=37\,500$; $C(28)=82\,000$. (b) No hay salto: los dos tramos dan $50\,000$ en $x=20$. (c) $50\,000+4\,000(x-20)=70\,000\Rightarrow x=\SI{25}{\cubic\metre}$.
\item (a) $3$ a la izquierda, estirar $\times2$, reflejar en el eje $x$, subir $1$. Dominio $[-3,\infty)$, rango $(-\infty,1]$; $(4,2)\mapsto(1,-3)$.
  (b) $f(x)=3$ si $x<3$; $f(x)=2x-3$ si $x\ge3$.
\item (a) $f^{-1}(x)=\dfrac{3x+1}{2-x}$, $x\ne2$. (b) $g^{-1}(x)=\sqrt{x}-1$, $x\ge0$.
\end{soluciones}
\end{respuestascortas}

\end{document}
